% ŠABLONA PRO PSANÍ ZÁVĚREČNÉ STUDIJNÍ PRÁCE
%%%%%%%%%%%%%%%%%%%%%%%%%%%%%%%%%%%%%%%%%%%%
% Autor: Jakub Dokulil (kubadokulil99@gmail.com)
% Tato šablona byla vytvořena tak, aby pomocí ní mohli v systému LaTeX soutěžící sázet své práce a zároveň odpovídala požadavkům na formátování vyplývajícím z wordové šablony umístěné na webu soc.cz.
%
\documentclass[12pt, a4paper,
oneside,      %% -- odkomentujte, pokud chcete svou práci mít pouze jednostrannou, mezera pro hřbet pak automaticky bude pouze na levé straně
%twoside,        %% -- pro oboustranné práce, mezera pro hřbet následně střídá strany.
%openright
]{report}

% --- odstraneni zbytkoveho textu "superiorSup" a pod. ---
%\AtBeginDocument{%
	% pojistka proti nechtenemu textu nactenemu z aux/toc
	%	\immediate\write16{(cleaning stray figureversions output...)}%
	%	\clearpage
	%	\thispagestyle{empty}
	% uplne vyprazdneni vseho, co by se objevilo mimo hlavni text
	%	\let\superiorSup\relax
	%	\let\textOsF\relax
	%	\let\textTOsF\relax
	%	\let\liningLF\relax
	%	\let\liningTLF\relax
	%	\let\tabularTab\relax
	%	\let\proportionalProp\relax
	%	\let\tabularmath\relax
	%	\let\proportionalmath\relax
	%	\let\fontspechyperref\relax
	% zajisteni, ze se nic nezobrazi pred titulni stranou
	%	\null
	%	\newpage
	%}
%% Nutné balíčky a nastavení
%%%%%%%%%%%%%%%%%%%%%%%%%%%%

%% Proměnné
\newcommand\obor{INFORMAČNÍ TECHNOLOGIE} %% -- napiš číslo a název tvého oboru
\newcommand\kodOboru{18-20-M/01} %% -- napiš číslo a název tvého oboru
\newcommand\zamereni{se zaměřením na počítačové sítě a programování} %% -- napiš číslo a název tvého oboru
\newcommand\skola{Střední škola průmyslová a umělecká, Opava} %% vyplň název školy
\newcommand\trida{IT4} %% vyplň jméno svého konzultanta
\newcommand\jmenoAutora{Filip Podeszwa}  %% vyplň své jméno
\newcommand\skolniRok{2025/26} %% vyplň rok
\newcommand\datumOdevzdani{7. 1. 2026} %% vyplň rok
\newcommand\nazevPrace{FreeVat} %% vyplň název své práce
\newcommand\popisPrace{Aplikace pro správu 3D modelů}

\title{\nazevPrace} %% -- Název tvé práce

\author{\jmenoAutora} %% -- tvé jméno
\date{\datumOdevzdani} %% -- rok, kdy píšeš SOČku

\usepackage[top=2.5cm, bottom=2.5cm, left=3.5cm, right=1.5cm]{geometry} %% nastaví okraje, left -- vnitřní okraj, right -- vnější okraj

\usepackage{float}
\usepackage[czech]{babel} %% balík babel pro sazbu v češtině
\usepackage[utf8]{inputenc} %% balíky pro kódování textu
\usepackage[T1]{fontenc}
\usepackage{cmap} %% balíček zajišťující, že vytvořené PDF bude prohledávatelné a kopírovatelné

\usepackage{graphicx} %% balík pro vkládání obrázků

\usepackage{subcaption} %% balíček pro vkládání podobrázků

\usepackage{hyperref} %% balíček, který v PDF vytváří odkazy

\captionsetup[subfigure]{labelformat=empty}

\linespread{1.25} %% řádkování
\setlength{\parskip}{0.5em} %% odsazení mezi odstavci


\usepackage[pagestyles]{titlesec} %% balíček pro úpravu stylu kapitol a sekcí
\titleformat{\chapter}[block]{\scshape\bfseries\LARGE}{\thechapter}{10pt}{\vspace{0pt}}[\vspace{-22pt}]
\titleformat{\section}[block]{\scshape\bfseries\Large}{\thesection}{10pt}{\vspace{0pt}}
\titleformat{\subsection}[block]{\bfseries\large}{\thesubsection}{10pt}{\vspace{0pt}}


\usepackage{tocloft} % Balíček umožní přizpůsobit vzhled tabulky obsahu
\setlength{\cftbeforechapskip}{0pt}  % Menší rozestup pro kapitoly
\setlength{\cftbeforesecskip}{0pt}   % Menší rozestup pro sekce

\setcounter{secnumdepth}{2}
\setcounter{tocdepth}{1}
\usepackage{fancyhdr}
\pagestyle{fancy}
\renewcommand{\headrulewidth}{0.025pt}

\usepackage{booktabs}

\usepackage{listings}
\usepackage{color}

\usepackage{url}

%% Balíčky co se můžou hodit :) 
%%%%%%%%%%%%%%%%%%%%%%%%%%%%%%%

\usepackage{pdfpages} %% Balíček umožňující vkládat stránky z PDF souborů, 

\usepackage{upgreek} %% Balíček pro sazbu stojatých řeckých písmen, třeba u jednotky mikrometr. Například stojaté mí: \upmu, stojaté pí: \uppi

\usepackage{amsmath}    %% Balíčky amsmath a amsfonts 
\usepackage{amsfonts}   %% pro sazbu matematických symbolů
\usepackage{esint}     %% pro sazbu různých integrálů (např \oiint)
\usepackage{mathrsfs}
\usepackage{helvet} % Helvet font
\usepackage{mathptmx} % Times New Roman
\usepackage{tikz}
\usetikzlibrary{positioning}
\makeatletter
\@namedef{ver@figureversions.sty}{9999/99/99}
\newcommand{\DeclareFigureVersion}[2]{}
\newcommand{\figureversion}[1]{}
\makeatother


\makeatletter
\providecommand{\superiorSup}{}
\providecommand{\textOsF}{}
\providecommand{\textTOsF}{}
\providecommand{\liningLF}{}
\providecommand{\liningTLF}{}
\providecommand{\tabularTab}{}
\providecommand{\proportionalProp}{}
\makeatother
\makeatletter
\providecommand{\superiorSup}{}
\providecommand{\textOsF}{}
\providecommand{\textTOsF}{}
\providecommand{\liningLF}{}
\providecommand{\liningTLF}{}
\providecommand{\tabularTab}{}
\providecommand{\proportionalProp}{}
\providecommand{\tabularmath}{}
\providecommand{\proportionalmath}{}
\makeatother

%\usepackage{Oswald} % Oswald font


%% makra pro sazbu matematiky
\newcommand{\dif}{\mathrm{d}} %% makro pro sazbu diferenciálu, místo toho
%% abych musel psát '\mathrm{d}' mi stačí napsat '\dif' což je mnohem 
%% kratší a mohu si tak usnadnit práci

\usepackage{listings}
\usepackage{xcolor}

\renewcommand{\lstlistingname}{Kód}% Listing -> Algorithm
\renewcommand{\lstlistlistingname}{Seznam programových kódů}% List of Listings -> List of Algorithms

%% Definice 
\lstdefinelanguage{JavaScript}{
	morekeywords=[1]{break, continue, delete, else, for, function, if, in,
		new, return, this, typeof, var, void, while, with},
	% Literals, primitive types, and reference types.
	morekeywords=[2]{false, null, true, boolean, number, undefined,
		Array, Boolean, Date, Math, Number, String, Object},
	% Built-ins.
	morekeywords=[3]{eval, parseInt, parseFloat, escape, unescape},
	sensitive,
	morecomment=[s]{/*}{*/},
	morecomment=[l]//,
	morecomment=[s]{/**}{*/}, % JavaDoc style comments
	morestring=[b]',
	morestring=[b]"
}[keywords, comments, strings]


\lstdefinelanguage[ECMAScript2015]{JavaScript}[]{JavaScript}{
	morekeywords=[1]{await, async, case, catch, class, const, default, do,
		enum, export, extends, finally, from, implements, import, instanceof,
		let, static, super, switch, throw, try},
	morestring=[b]` % Interpolation strings.
}

\lstalias[]{ES6}[ECMAScript2015]{JavaScript}

% Nastavení barev
% Requires package: color.
\definecolor{mediumgray}{rgb}{0.3, 0.4, 0.4}
\definecolor{mediumblue}{rgb}{0.0, 0.0, 0.8}
\definecolor{forestgreen}{rgb}{0.13, 0.55, 0.13}
\definecolor{darkviolet}{rgb}{0.58, 0.0, 0.83}
\definecolor{royalblue}{rgb}{0.25, 0.41, 0.88}
\definecolor{crimson}{rgb}{0.86, 0.8, 0.24}

% Nastavení pro Python
\lstdefinestyle{Python}{
	language=Python,
	backgroundcolor=\color{white},
	basicstyle=\ttfamily,
	breakatwhitespace=false,
	breaklines=false,
	captionpos=b,
	columns=fullflexible,
	commentstyle=\color{mediumgray}\upshape,
	emph={},
	emphstyle=\color{crimson},
	extendedchars=true,  % requires inputenc
	fontadjust=true,
	frame=single,
	identifierstyle=\color{black},
	keepspaces=true,
	keywordstyle=\color{mediumblue},
	keywordstyle={[2]\color{darkviolet}},
	keywordstyle={[3]\color{royalblue}},
	literate=%
	{á}{{\'a}}1 {č}{{\v{c}}}1 {ď}{{\v{d}}}1 {é}{{\'e}}1 {ě}{{\v{e}}}1
	{í}{{\'i}}1 {ň}{{\v{n}}}1 {ó}{{\'o}}1 {ř}{{\v{r}}}1 {š}{{\v{s}}}1
	{ť}{{\v{t}}}1 {ú}{{\'u}}1 {ů}{{\r{u}}}1 {ý}{{\'y}}1 {ž}{{\v{z}}}1,		
	numbers=left,
	numbersep=5pt,
	numberstyle=\tiny\color{black},
	rulecolor=\color{black},
	showlines=true,
	showspaces=false,
	showstringspaces=false,
	showtabs=false,
	stringstyle=\color{forestgreen},
	tabsize=2,
	title=\lstname,
	upquote=true  % requires textcomp	
}


\lstdefinestyle{JSES6Base}{
	backgroundcolor=\color{white},
	basicstyle=\ttfamily,
	breakatwhitespace=false,
	breaklines=false,
	captionpos=b,
	columns=fullflexible,
	commentstyle=\color{mediumgray}\upshape,
	emph={},
	emphstyle=\color{crimson},
	extendedchars=true,  % requires inputenc
	fontadjust=true,
	frame=single,
	identifierstyle=\color{black},
	keepspaces=true,
	keywordstyle=\color{mediumblue},
	keywordstyle={[2]\color{darkviolet}},
	keywordstyle={[3]\color{royalblue}},
	literate=%
	{á}{{\'a}}1 {č}{{\v{c}}}1 {ď}{{\v{d}}}1 {é}{{\'e}}1 {ě}{{\v{e}}}1
	{í}{{\'i}}1 {ň}{{\v{n}}}1 {ó}{{\'o}}1 {ř}{{\v{r}}}1 {š}{{\v{s}}}1
	{ť}{{\v{t}}}1 {ú}{{\'u}}1 {ů}{{\r{u}}}1 {ý}{{\'y}}1 {ž}{{\v{z}}}1,		
	numbers=left,
	numbersep=5pt,
	numberstyle=\tiny\color{black},
	rulecolor=\color{black},
	showlines=true,
	showspaces=false,
	showstringspaces=false,
	showtabs=false,
	stringstyle=\color{forestgreen},
	tabsize=2,
	title=\lstname,
	upquote=true  % requires textcomp
}

\lstdefinestyle{JavaScript}{
	language=JavaScript,
	style=JSES6Base,
}
\lstdefinestyle{ES6}{
	language=ES6,
	style=JSES6Base
}

\setlength{\headheight}{15pt}

\usepackage{fancyhdr}
\setlength{\headheight}{15pt}

% ===============================
% Záhlaví a zápatí
% ===============================
\fancyhf{}                                % vymazání výchozího obsahu
\fancyhead[L]{Závěrečná práce}
\fancyhead[C]{Filip Podeszwa IT4}
\fancyhead[R]{2025/2026}
\fancyfoot[C]{\thepage}

\renewcommand{\headrulewidth}{0.4pt}
\renewcommand{\footrulewidth}{0pt}

% Styl pro kapitoly (\chapter)
\fancypagestyle{plain}{
	\fancyhf{}
	\fancyhead[L]{Závěrečná práce}
	\fancyhead[C]{Filip Podeszwa IT4}
	\fancyhead[R]{2025/2026}
	\fancyfoot[C]{\thepage}
	\renewcommand{\headrulewidth}{0.4pt}
}




%% Začátek dokumentu
%%%%%%%%%%%%%%%%%%%%
\begin{document}
	
	\pagestyle{empty}
	\pagenumbering{Roman}
	\pagenumbering{gobble} % žádné číslování
	%\clearpage
	
	%% Titulní stránka s informacemi
	%%%%%%%%%%%%%%%%%%%%%%%%%%%%%%%%%%%%%%%%
	
	{\fontfamily{phv}\selectfont
		%% Logo školy
		\begin{figure}[h]
			\centering
			\includegraphics[width=0.6\linewidth]{image/logo-skoly.png} 
		\end{figure}
		
		
		%% Hlavička práce a její název (viz proměnná \nazev prace)
		%% \sffamily %%% bezpatkové písmo - sans serif
		{\bfseries %%% písmo na stránce je tučně
			\begin{center}
				\vspace{0.025 \textheight}
				\LARGE{ZÁVĚREČNÁ STUDIJNÍ PRÁCE}\\
				\large{dokumentace}\\
				\vspace{0.075 \textheight}
				\LARGE {\nazevPrace}\\
				\large {\popisPrace}\\
			\end{center}  
		}%%%
		
		\begin{figure}[h]
			\centering
			\includegraphics[width=0.4\linewidth]{img/favicon.png} 
		\end{figure}
		
		\vspace{0.02 \textheight}
		\begin{table}[h!]
			\begin{tabular}{ll}
				\textbf{Autor:} & \jmenoAutora\\ 
				\textbf{Obor:} & \kodOboru { } \obor\\
				\textbf{} & \zamereni\\
				\textbf{Třída:} & \trida\\
				\textbf{Školní rok:} & \skolniRok\\
			\end{tabular}
			
		\end{table}		
	}
	
	\clearpage %% Zalomení dvojstránky
	
	%% Stránka obsahující poděkování a prohlášení
	%%%%%%%%%%%%%%%%%%%%%%%%%%%%%%%%%%%%%%%%%%%%%%%%%%%%%%%%
	
	%% Poděkování - nepovinné
	%%%%%%%%%%%%%%%%%%%%%%%%%%%%
	
	\noindent{\large{\bfseries{Poděkování}\\}}
	\noindent Rád bych poděkoval všem umělým inteligencím, zejména \textbf{Google Gemini}, která mi pomáhala od začátku do konce, dále \textbf{ChatGPT} a taky \textbf{DeepSeek}, který není vůbec špatný. Zpočátku mi pomáhal dokonce i \textbf{GitHub Copilot}, ale ten se na mě téměř před čtvrt rokem vykašlal :(
	
	\noindent Co se týče lidí, rád bych poděkoval panu \textbf{Ing. Petru Grussmannovi} za pomoc v Kubernetesu, která se ukázala být k ničemu, ale to je můj problém, protože jsem aplikaci již nestihl zdockerovat. Nasazení aplikace na web se budu věnovat někdy později.
	
	\vspace*{0.38\textheight} %% Vertikální mezeru je možné upravit
	
	%% Prohlášení - povinné
	%%%%%%%%%%%%%%%%%%%%%%%%%%%%
	\noindent{\large{\bfseries{Prohlášení}\\}}  %% uprav si koncovky podle toho na jaký rod se cítíš, vypadá to pak lépe :) 
	\noindent{Prohlašuji, že jsem závěrečnou práci vypracoval samostatně a uvedl veškeré použité 
		informační zdroje.\\}
	\noindent{Souhlasím, aby tato studijní práce byla použita k výukovým a prezentačním účelům na Střední průmyslové a umělecké škole v Opavě, Praskova 399/8.}
	\vfill
	\noindent{V Opavě \datumOdevzdani\\}
	\noindent
	\begin{minipage}{\linewidth}
		\hspace{9.5cm} 
		\begin{tabular}{@{}p{6cm}@{}}
			\dotfill \\
			Podpis autora
		\end{tabular}
	\end{minipage}
	
	\clearpage
	
	
	%% Stránka obsahující abstrakt (anotaci)
	%%%%%%%%%%%%%%%%%%%%%%%%%%%%%%%%%%%%%%%%%%%%%%%%%%%%%%%%	
	
	%% Abstrakt v češtině
	%%%%%%%%%%%%%%%%%%%%%%%%%%%%
	\noindent{\Large{\bfseries{Abstrakt}}}
	
	\noindent\textbf{FreeVat} je webová platforma zaměřená na komunitní \textbf{sdílení}, \textbf{prohlížení} a \textbf{správu} 3D modelů. Hlavním cílem aplikace je poskytnout uživatelům intuitivní prostředí pro prezentaci jejich digitální tvorby \textbf{bez nutnosti instalace} specializovaného softwaru. Platforma slouží jako centrální úložiště, které propojuje tvůrce 3D obsahu s koncovými uživateli, kteří tyto modely mohou využít ve \textbf{vlastních projektech, hrách nebo při 3D tisku}. 
	
	\noindent Aplikace umožňuje uživatelům \textbf{nahrávat modely, zobrazovat je ve 360°, spravovat je a taky mazat.} Přihlašování je možné jak klasickou metodou jméno/email/heslo, tak pomocí \textbf{Google} či \textbf{GitHubu}. Poskytuje taky prostředí pro psaní komentářů na ostatní modely a zobrazuje základní informace o nich.
	
	\vspace{18pt}
	
	\noindent{\large{\bfseries{Klíčová slova}}}
	
	\noindent FreeVat, 3D grafika, 3D modely, Django, Python, Three.js, HTML, CSS, JavaScript, Tailwind, Adobe Photoshop, PostgreSQL, Node.js, open-source, správa digitálního obsahu, webová aplikace, interaktivní prohlížeč, web, webové stránky, digitální grafika, programování, modelování
	
	\vspace{1.5cm}
	
	%% Abstrakt v angličtině
	%%%%%%%%%%%%%%%%%%%%%%%%%%%%	
	\noindent{\Large{\bfseries{Abstract}}}
	
	\noindent\textbf{FreeVat} is a web platform focused on community \textbf{sharing}, \textbf{viewing}, and \textbf{managing} 3D models. The main goal of the application is to provide users with an intuitive environment for presenting their digital creations \textbf{without the need to install} specialized software. The platform serves as a central repository that connects 3D content creators with end users who can use these models in \textbf{their own projects, games, or 3D printing}. 
	
	\noindent The application allows users to \textbf{upload models, view them in 360°, manage them, and delete them}. Users can log in using the classic name/email/password method, but it is also possible via \textbf{Google} or \textbf{GitHub}. It also provides an environment for commenting on other models and displays basic information about them.
	
	\vspace{18pt}
	
	\noindent{\large{\bfseries{Keywords}}}
	
	\noindent FreeVat, 3D graphics, 3D models, Django, Python, Three.js, HTML, CSS, JavaScript, Tailwind, Adobe Photoshop, PostgreSQL, Node.js, open-source, digital content management, web app, interactive viewer, web, web pages, digital graphics, programming, modeling
	
	
	
	\clearpage %% Zalomení stránky
	
	%% Stránka s generovaným obsahem
	%%%%%%%%%%%%%%%%%%%%%%%%%%%%%%%%%%%%%%%	
	
	\tableofcontents %% Vygeneruje tabulku s obsahem
	
	\clearpage
	
	\pagenumbering{arabic}
	\setcounter{page}{1}
	\pagestyle{fancy}
	
	%% Stránka s úvodem - povinná část
	%%%%%%%%%%%%%%%%%%%%%%%%%%%%%%%%%%%%%%%		
	\chapter*{Úvod}
	%Tento příkaz vytvoří novou kapitolu s názvem "Úvod" ve vašem dokumentu.
	%Hvězdička * u příkazu \chapter* znamená, že tato kapitola nebude mít číslo. Ve výsledném dokumentu se tedy objeví jako "Úvod" bez předcházejícího čísla kapitoly, které se obvykle zobrazuje u číslovaných kapitol.
	%Tento příkaz také znamená, že kapitola se automaticky neobjeví v obsahu, protože LaTeX standardně zahrnuje do obsahu pouze číslované kapitoly.
	\addcontentsline{toc}{chapter}{Úvod}
	%Tento příkaz ručně přidává záznam do obsahu.
	%První parametr toc označuje, že přidáváme záznam do Table of Contents (obsahu).
	%Druhý parametr chapter specifikuje úroveň záznamu. V tomto případě říkáme, že přidávaný záznam má být považován za kapitolu.
	%Třetí parametr Úvod je text, který se objeví v obsahu. V tomto případě bude v obsahu zobrazen název "Úvod".	
	
	Dřív, než se ponoříme do samotného popisu vývoje aplikace, chtěl bych uvést na pravou míru, že toto byl můj první "velký" projekt, co se IT týče. Nikdy předtím jsem nic takového nedělal. Toto bylo poprvé, co jsem musel pracovat s více věcmi najednou. Nestačilo např. jen něco vytvořit v jazyce \textit{C} či \textit{HTML}, musel jsem do sebe spojit více technologií a zajistit, aby spolu řádně komunikovaly.
	
	\noindent Proto prosím berte na vědomí, že postupy v mé aplikaci mnohdy nemusí zcela vyhovovat zažitým praktikám, občas se i vyskytuje malinké množství duplicitního kódu, který samozřejmě v žádné správně naprogramované aplikaci není úplně žádoucí. Stejně tak nemám do hloubky prostudované frameworky, které jsem k vývoji použil. Je tudíž možné, že udělám nějaký překlep nebo i něco špatně interpretuji.
	
	%Tipy k psaní úvodu
	%Je povinný, nadpis neměňte, rozsah - max. 1 strana. 
	%Tato část práce obsahuje: 
	%* náhled do řešené problematiky, zdůvodnění volby problematiky, 
	%* předem definované cíle práce, 
	%* motivaci pro další čtení textu včetně stručného uvedení obsahu následujících kapitol 
	
	
	\chapter{Využité technologie}
	
	Tato kapitola obsahuje podrobný popis technologií, které byly k projektu \textbf{FreeVat} využity. Jsou to jak technologie, co se starají o vzhled, tak i ty, co jsou zodpovědné za vnitřní logiku i běh aplikace. Všechny zdroje jsou uvedeny na konci dokumentace v sekci \textbf{Zdroje}.
	
	\section{Front-end}
	
	V této sekci jsou popsány technologie využité při vývoji front-endové části aplikace, která tvoří rozhraní pro interakci s uživatelem. Základním stavebním kamenem je kombinace jazyků \textbf{HTML, CSS a JavaScriptu}, doplněná o framework \textbf{Tailwind CSS} pro efektivní tvorbu moderního a responzivního designu. Klíčovou roli zde hraje knihovna \textbf{Three.js}, která do prostředí webového prohlížeče integruje pokročilou 3D grafiku, zatímco o vizuální prvky a grafické podklady se postaral software \textbf{Adobe Photoshop}. Cílem této části je objasnit výběr těchto nástrojů a jejich přínos pro výslednou vizuální podobu a funkčnost uživatelského prostředí.
	
	\vspace*{1.5cm}
	
	\subsection{Tailwind CSS}
	
	Při stylování stránek jsem se rozhodl využít framework Tailwind CSS. Je lepší než klasické CSS, protože využívá tzv. \textbf{utility-first}, což jsou předpřipravené třídy, které lze spojovat a ušetřit tak veliké množství kódu.
	
	\begin{figure}[H] % [H] se snaží udržet obrázek PŘESNĚ ZDE. Pokud to nejde, zkus [hbt!]
		\centering % Vycentruje celý blok na stránce
		
		% --- PRVNÍ OBRÁZEK VLEVO ---
		\begin{subfigure}[b]{0.48\textwidth} % Šířka bloku je 48% stránky
			\centering
			% ZDE NAHRAĎ CESTU K SOUBORU:
			\includegraphics[width=\textwidth]{img/css_example.png}
			\caption{Klasické CSS styly} % Malý popisek (a)
			\label{fig:levy-obrazek}
		\end{subfigure}
		\hfill % Pružná mezera mezi obrázky
		% --- DRUHÝ OBRÁZEK VPRAVO ---
		\begin{subfigure}[b]{0.48\textwidth} % Šířka bloku je 48% stránky
			\centering
			% ZDE NAHRAĎ CESTU K SOUBORU:
			\includegraphics[width=\textwidth]{img/tailwind_example.png}
			\caption{Tailwind "utility-first"} % Malý popisek (b)
			\label{fig:pravy-obrazek}
		\end{subfigure}
	\end{figure}
	%%%%% KONEC BLOKU S OBRÁZKY %%%%%
	
	\clearpage
	
	\subsection{Three.js}
	
	Klíčovou roli v aplikaci hraje Three.js. Je to populární open-source knihovna, která výrazně zjednodušuje tvorbu interaktivní 3D grafiky přímo v prohlížeči. Funguje jako nadstavba nad rozhraním WebGL, což umožňuje plně využít výkon grafické karty pro hardwarově akcelerované renderování v reálném čase.
	
	\begin{figure}[H]
		\includegraphics[width=\textwidth]{img/threejs_code.png}
		\caption*{Ukázka kódu s použitím Three.js}
	\end{figure}
	
	\vspace*{1cm}
	
	\subsection{Adobe Photoshop 2026}
	
	Ačkoliv je Photoshop primárně rastrový editor, je možné v něm tvořit i nějaké to jednoduché \textbf{SVG}. V tomto programu byly vytvořeny \textbf{vlajky a loga aplikace}.
	
	\clearpage
	
	\section{Back-end}
	
	Tato sekce se obsahuje popis technologií při návrhu a implementaci back-endové části aplikace, která zajišťuje logiku systému a správu dat. Hlavním pilířem je framework Django - robustní základ pro serverovou část, doplněný o nástroje Crispy Forms pro pokročilé zpracování formulářů a Rosetta pro zajištění vícejazyčného rozhraní. Pro bezpečné a výkonné ukládání informací byl zvolen databázový systém PostgreSQL, zatímco integraci externích služeb a bezpečné přihlašování zajišťuje protokol OAuth. Celý vývojový proces a efektivní sestavení aplikace (buildování) pak doplňuje prostředí Node.js spolu s moderním nástrojem Vite.
	
	\vspace*{1cm}
	
	\subsection{Django}
	
	Jako základ projektu jsem zvolil framework Django, postavený na čitelném jazyku Python. Díky přístupu „batteries-included“ a integrovanému ORM výrazně urychluje vývoj a zajišťuje vysokou úroveň bezpečnosti proti běžným útokům. Architektura je navržena pro snadnou škálovatelnost a podporuje čistý, pragmatický design celého systému.
	
	\begin{figure}[H]
		\centering
		\includegraphics[width=12cm]{img/django.png}
	\end{figure}
	
	\clearpage
	
	\subsection{Python knihovny}
	
	\begin{figure}[H]
		\centering
		\includegraphics[width=14cm]{img/python_wide.png}
	\end{figure}
	
	\subsubsection{Rosetta}
	Protože jsem chtěl udělat web ve 3 různých jazycích, potřeboval jsem efektivně vytvářet a spravovat překlady, a na to jsem využil knihovnu Rosetta. Díky této knihovně je možné spravovat překlady přímo v admin panelu jednoduše a přehledně.
	
	\vspace*{1.2cm}
	
	\subsubsection{Crispy Forms}
	Tento balík jsem využil na tvorbu a stylování uživatelských formulářů. Návrh tak šel mnohem rychleji, než kdybych je měl programovat od píky. Za zmínku taky stojí, že pomocí malé knihovničky je možné Crispy Forms propojit s Tailwindem a je možné si tak vytvářet vlastní stylované formuláře, než používat ošklivý vestavěný design.
	
	\vspace*{1.2cm}
	
	\subsubsection{PyYAML}
	Díky této knihovně jsem databázové dumpy mohl uložit do formátu \textbf{YAML}, který je pro člověka mnohem lépe čitelný než obyčejný \textbf{JSON} a za mě je i spolehlivější.
	
	\clearpage
	
	\subsection{PostgreSQL}
	
	Django sice má v sobě vestavěnou databázi \textbf{SQLite}, ale ta je hodně omezená a pro větší projekty nevhodná. Proto jsem sáhnul po robustnějším řešení a nainstaloval \textbf{PostgreSQL}, což je velice výkonný \textbf{objektově-relační} databázový systém. Je známý svojí spolehlivostí a integritou dat. A hlavně je open-source, takže zadarmo!
	
	\begin{figure}[H]
		\centering
		\includegraphics[width=3cm]{img/postgresql.png}
	\end{figure}
	
	\subsection{OAuth}
	
	\noindent % Zabrání odsazení prvního řádku
	\begin{minipage}[c]{0.68\textwidth} % Text zabere 75 % šířky
		OAuth je otevřený protokol, který v aplikaci zajišťuje bezpečné ověřování uživatelů prostřednictvím třetích stran, jako je Google, Facebook nebo GitHub. Díky tomu se uživatelé nemusí starat o hesla a jednoduše se kdykoli přihlásí a odhlásí jedním kliknutím.
	\end{minipage}
	\hfill % Pružná mezera mezi textem a logem
	\begin{minipage}[c]{0.27\textwidth} % Logo zabere 20 % šířky
		\centering
		\includegraphics[width=5cm]{img/oauth2.png} % Nahraď svou cestou k logu
	\end{minipage}
	
	\begin{figure}[H]
		\centering
		\includegraphics[width=7.4cm]{img/login.png}
		\caption*{Přihlašovací formulář}
	\end{figure}
	
	\clearpage
	
	\subsection{Node.js a Vite}
	Pro integraci technologií \textbf{Tailwind CSS} a \textbf{Three.js} bylo nezbytné zvolit moderní vývojové prostředí schopné efektivně kompilovat zdrojové kódy a styly. Klíčovou roli zde hraje \textbf{Node.js}, který poskytuje runtime prostředí pro buildovací procesy. 
	\noindent Nástroj \textbf{Vite} zajišťuje extrémně rychlou obsluhu kódu a jeho optimalizaci pro produkční nasazení. Díky technologii Hot Module Replacement (HMR) umožňuje Vite promítat změny v JavaScriptu a CSS okamžitě do prohlížeče, což zásadně zrychluje a zpříjemňuje celý proces vývoje aplikace.
	
	\begin{figure}[H] % [H] se snaží udržet obrázek PŘESNĚ ZDE. Pokud to nejde, zkus [hbt!]
		\centering % Vycentruje celý blok na stránce
		
		% --- PRVNÍ OBRÁZEK VLEVO ---
		\begin{subfigure}[b]{0.48\textwidth} % Šířka bloku je 48% stránky
			\centering
			% ZDE NAHRAĎ CESTU K SOUBORU:
			\includegraphics[width=6.5cm]{img/nodejs.png}
		\end{subfigure}
		\hfill % Pružná mezera mezi obrázky
		% --- DRUHÝ OBRÁZEK VPRAVO ---
		\begin{subfigure}[b]{0.48\textwidth} % Šířka bloku je 48% stránky
			\centering
			% ZDE NAHRAĎ CESTU K SOUBORU:
			\includegraphics[width=5cm]{img/vite.png}
		\end{subfigure}
	\end{figure}

		
	\chapter{Vývoj}
	V této kapitole jsou podrobně rozepsány všechny verze projektu. V každé je uvedeno, co bylo přidáno, co odebráno nebo co opraveno.
	
	\clearpage
	
	\section{Indev 0.0.1}
	
	V červnu létě 2025 jsem udělal ten úplně první krok - vytvořil jsem repozitář na GitHubu. Hned nato jsem přidal soubor \textbf{README.md}, do kterého jsem vypsal všechny informace o projektu - popis, závislosti, cíle...
	
	\begin{figure}[H]
		\centering
		\includegraphics[width=\textwidth]{img/readme.png}
	\end{figure}
	
	\noindent Vzápětí jsem zapnul \textbf{PyCharm}, zprovoznil \textbf{virtuální prostředí} a nainstaloval Django aplikaci.
	\noindent Během toho jsem přemýšlel i nad jménem. Napadlo mě více možností, ale vyhrál \textbf{FreeVat}.
	\noindent Nakonec jsem nainstaloval pár nezbytných závislostí.
	
	\clearpage
	
	\section{Indev 0.0.5}
	
	Původně jsem místo \textbf{Tailwindu} používal \textbf{Bootstrap}, ale později jsem zjistil, že je omezenější a není tak flexibilní jako \textbf{Tailwind}.
	
	\begin{figure}[H]
		\centering
		\includegraphics[width=\textwidth]{img/bootstrap_tailwind.png}
	\end{figure}
	
	\noindent Touto dobou už jsem pracoval na záhlaví, navigačním menu a zápatí webu. Dokonce se mi podařilo zprovoznit tlačítko pro přepínání jazyků, a tak už tehdy byl web v Češtině, Angličtině a Němčině. Přepínání je zprovozněno pomocí JavaScriptu.
	
	\vspace*{1cm}
	
	\begin{figure}[H]
		\centering
		
		\begin{subfigure}[b]{0.28\textwidth}
			\centering
			\includegraphics[width=\textwidth]{img/cz.png}
		\end{subfigure}% <--- TENTO ZNAK JE DŮLEŽITÝ (odstraní mezeru)
		\hfill 
		\begin{subfigure}[b]{0.28\textwidth}
			\centering
			\includegraphics[width=\textwidth]{img/gb.png}
		\end{subfigure}%
		\hfill
		\begin{subfigure}[b]{0.28\textwidth}
			\centering
			\includegraphics[width=\textwidth]{img/de.png}
		\end{subfigure}
	\end{figure}
	
	\clearpage
	
	\section{Indev 0.0.7}
	
	V říjnu jsem dokončil záhlaví, navigaci a zápatí. Přidal jsem dokonce vyhledávací lištu (která ovšem zatím není funkční :)
	
	\begin{figure}[H]{\textwidth}
		\centering
		\includegraphics[width=\textwidth]{img/navbar.png}
	\end{figure}
	
	\begin{figure}[H]{\textwidth}
		\centering
		\includegraphics[width=\textwidth]{img/footer.png}
	\end{figure}
	
	\vspace*{1.2cm}
	
	\noindent Dokonce v té době proběhl první pokus o implementaci 3D prohlížeče na web a dokonce se to i podařilo, ačkoliv to nefungovalo tak, jak by mělo.
	
	\begin{figure}[H]{\textwidth}
		\centering
		\includegraphics[width=\textwidth]{img/npm_run_dev.png}
	\end{figure}
	
	\clearpage
	
	
	\section{Indev 0.0.9}
	
	Zprovoznění formulářů pro registraci nových účtů a přihlašování se ke stávajícím účtům. Funguje jak klasické přihlášení, tak i OAuth přes Google a GitHub.
	
	\begin{figure}[H]{\textwidth}
		\centering
		\includegraphics[width=\textwidth]{img/profile.png}
	\end{figure}
	
	\pagebreak
	
	\section{Pre-Alpha 0.1.0}
	
	\noindent Tohle byl docela velký update, protože vyšla první Pre-Alpha. Podařilo se mi zprovoznit formulář pro nahrávání 3D modelů, které se nyní krásně ukládají do databáze i se všemi daty.
	
	\begin{figure}[H]{\textwidth}
		\centering
		\includegraphics[width=\textwidth]{img/upload2.png}
	\end{figure}
	
	\clearpage
	
	\noindent Taky 3D viewer se mi podařilo plně zprovoznit. Nyní je možné model posouvat, přibližovat a rotovat.
	
	\begin{figure}[H]{\textwidth}
		\centering
		\includegraphics[width=\textwidth]{img/3d_viewer.png}
	\end{figure}
	
	\clearpage
	
	\section{Alpha 1.0.0}
	
	A je to tady! Vyšla historicky první Alpha verze mého projektu. V této fázi je připravená k používání - soubory se správně nahrávají, 3D zobrazení funguje, uživatelé můžou stahovat modely a psát komentáře
	
	
	\noindent Domovská obrazovka je hezky nadesignovaná
	
	\begin{figure}[H]{\textwidth}
		\centering
		\includegraphics[width=\textwidth]{img/home.png}
	\end{figure}
	
	
	\noindent Uživatelé mohou upravovat nebo mazat své nahrané modely
	
	\begin{figure}[H]{\textwidth}
		\centering
		\includegraphics[width=\textwidth]{img/user_models.png}
	\end{figure}
	
	\clearpage
	
	\noindent O modelu se zobrazují technické informace, jako je počet polygonů a vrcholů, velikost souboru a formát modelu. Zobrazují se i čtyři náhledové obrázky.
	
	\begin{figure}[H]{\textwidth}
		\centering
		\includegraphics[width=\textwidth]{img/model_detail.png}
	\end{figure}
	
	\clearpage
	
	\chapter{Splněné a nesplněné cíle}
	
	\vspace*{0.5cm}
	
	\section{Splněno}
	
	\begin{itemize}
		\item ukládání modelů do databáze i s jejich daty
		\item provádění základních CRUD operací s modely
		\item prohlížení modelů ve 360°
		\item posouvání, přibližování a otáčení modelů
		\item registrace, přihlašování a mazání uživatelských účtů
		\item zobrazení základních informací o modelech
		\item vytvoření vícejazyčné webové aplikace
		\item řazení modelů podle názvu, data nahrání a velikosti
	\end{itemize}
	
	\vspace*{0.3cm}
	
	\section{Nesplněno/Neimplementováno}
	
	\begin{itemize}
		\item podpora formátů 3D softwarů
		\item plně responzivní design na PC, tabletu i mobilu
		\item vyhledávání 3D modelů podle názvu
		\item zdockerování aplikace
		\item a samozřejmě mnoho dalšího
	\end{itemize}
	
	
	\chapter*{Závěr}
	\addcontentsline{toc}{chapter}{Závěr}
	
	Hlavním cílem tohoto projektu bylo vytvořit komplexní webovou aplikaci schopnou ukládat a interaktivně zobrazovat 3D modely různých formátů. Realizace tohoto záměru mi umožnila nejen dosáhnout vytyčených technických cílů, ale především získat hluboké praktické porozumění frameworku \textbf{Django} a osvojit si moderní přístupy k tvorbě uživatelského rozhraní pomocí \textbf{Tailwind CSS}, kterému jsem díky tomuto projektu porozuměl do hloubky. 
	
	\vspace*{0.2cm}
	
	\noindent Jelikož se jednalo o můj první velký projekt tohoto rozsahu, musel jsem čelit mnoha výzvám, od návrhu databázové struktury až po integraci 3D grafiky v reálném čase. Celý proces vývoje mě naučil efektivně řešit komplexní problémy a propojovat různorodé technologie do jednoho stabilního celku. Aplikace FreeVat je tak výsledkem dlouhodobého učení a představuje funkční základ, který jsem připraven v budoucnu dále rozšiřovat o nové moduly a optimalizace.
	
	
	\begin{thebibliography}{99}
		
		\bibitem{htmlSpec}
		W3C. \textit{HTML Living Standard}. [online]. Dostupné z: \url{https://html.spec.whatwg.org/} [cit. 2025-01-05].
		
		\bibitem{cssSpec}
		W3C. \textit{CSS Specifications}. [online]. Dostupné z: \url{https://www.w3.org/Style/CSS/Overview.en.html} [cit. 2025-01-05].
		
		\bibitem{javascriptMDN}
		Mozilla Developer Network. \textit{JavaScript Guide}. [online]. Dostupné z: \url{https://developer.mozilla.org/en-US/docs/Web/JavaScript} [cit. 2025-01-06].
		
		\bibitem{threejsDocs}
		Three.js Authors. \textit{Three.js Documentation}. [online]. Dostupné z: \url{https://threejs.org/docs/} [cit. 2025-01-06].
		
		\bibitem{tailwindDocs}
		Tailwind Labs. \textit{Tailwind CSS Documentation}. [online]. Dostupné z: \url{https://tailwindcss.com/docs} [cit. 2025-01-07].
		
		\bibitem{photoshopDocs}
		Adobe. \textit{Adobe Photoshop Documentation}. [online]. Dostupné z: \url{https://helpx.adobe.com/photoshop/user-guide.html} [cit. 2025-01-07].
		
		\bibitem{djangoDocs}
		Django Software Foundation. \textit{Django Documentation}. [online]. Dostupné z: \url{https://docs.djangoproject.com/} [cit. 2025-01-08].
		
		\bibitem{djangoI18n}
		Django Software Foundation. \textit{Internationalization and Localization}. [online]. Dostupné z: \url{https://docs.djangoproject.com/en/stable/topics/i18n/} [cit. 2025-01-08].
		
		\bibitem{crispyForms}
		Django Crispy Forms. \textit{Crispy Forms Documentation}. [online]. Dostupné z: \url{https://django-crispy-forms.readthedocs.io/} [cit. 2025-01-09].
		
		\bibitem{postgresDocs}
		PostgreSQL Global Development Group. \textit{PostgreSQL Documentation}. [online]. Dostupné z: \url{https://www.postgresql.org/docs/} [cit. 2025-01-09].
		
		\bibitem{oauthSpec}
		IETF. \textit{The OAuth 2.0 Authorization Framework}. [online]. Dostupné z: \url{https://datatracker.ietf.org/doc/html/rfc6749} [cit. 2025-01-10].
		
		\bibitem{nodejsDocs}
		OpenJS Foundation. \textit{Node.js Documentation}. [online]. Dostupné z: \url{https://nodejs.org/en/docs} [cit. 2025-01-10].
		
		\bibitem{viteDocs}
		Vite Team. \textit{Vite Documentation}. [online]. Dostupné z: \url{https://vitejs.dev/guide/} [cit. 2025-01-11].
		
		
	\end{thebibliography}
	
	%% tabulky
	%\listoftables
	
	\appendix %% začínají přílohy
	
	\titleformat{\chapter}[block]{\scshape\bfseries\LARGE}{Příloha \thechapter}{10pt}{\vspace{0pt}}[\vspace{-22pt}] %% nastavení nadpisu u příloh
	
	
	
\end{document}